% s-342-p342-m002; original PDF 342
\Aphorism{s-342-p342-m002}{393}
Man verdirbt sich die transcendenten Probleme durch Einmischung der empirischen Thatsache Tod. Beispielsweise: warum müssten wir aus diesem Leben geschieden sein, um in ein anderes übergehen zu können? Haben wir an einer Mehrheit von Welten Theil, so besteht diese Theilhaberschaft ganz unabhängig davon, wie weit wir im zeitlichen Verlauf einer dieser Welten gelangt sind. Wiedergeburt bedeutet nur ein Sein in mehr als einem Ich, sagt aber nichts über das empirische Zugleich- oder Nacheinanderleben dieser verschiedenen Subjecte. — Oder warum müssten wir gestorben sein, um uns als Wesen eines vierdimensionalen Raumes zu bethätigen, gesetzt, dass wir überhaupt über den dreifach ausgedehnten Raum hinaus organisirt sind? — Bedenken hinsichtlich des Todes führten zur Philosophie, und schliesslich hört gerade der Tod auf, ein philosophisches Problem zu sein!
\par

